\subsection{Strict harmonic comparison on the constructed core}\label{subsec:fc-line-separated-harmonic}

\begin{theorem}[Strict harmonic boundary bound on the core]\label{thm:fc-separated-harmonic-strict-bound}
\lean{EnvelopeProperty.Berkovich.RelativeSpectrum.FCStandardAffinoid.IsHarmonic.lt_of_boundary_le_of_exists_lt}
\leanok
\uses{lem:fc-separated-core-connected-boundary, thm:fc-harmonic-component-maximum,
  thm:fc-harmonic-weak-maximum}
Let $U$ have a separated presentation and nonempty strict core. If an actual
harmonic function $H$ has boundary values at most $R$, with one actual
boundary value strictly smaller than $R$, then $H(x)<R$ for every point
$x$ of its strict core.
\end{theorem}

\begin{proof}
The actual weak maximum principle gives $H\le R$ on the whole carrier.
If equality holds at a strict-core point, the actual componentwise harmonic
maximum theorem makes $H$ constant on the interior component containing
that point. The connected strict core is contained in this component.
Every actual boundary point lies in the closure of the core, so continuity
makes every boundary value equal to $R$, a contradiction.
The separated presentation will always be obtained by construction in the
subharmonic maximum-principle application.
\end{proof}
